\BeleskeID{04-s006}
% element: 04-s006:e01 — Primer 326.23 u osnovi 7 i sva težinska množenja; 167+17/49
% element: 04-s006:e02 — Oba decimalna niza iz originala za 17/49 i kompletnu vrednost
% element: 04-s006:e03 — Zaključak o konačnom ili beskonačnom zapisu pri promeni osnove
% element: 04-s006:e04 — Komentar o školskom računanju razlomcima i želji za standardnim zapisom
% element: 04-s006:d01

Zapis $167+17/49$ daje tačnu vrednost, kao u školskom računanju razlomcima. Kada tražimo standardan pozicioni zapis, želimo cifre pre i posle decimalne tačke. Izraz da razlomci „ne postoje“ odnosi se na taj zahtev za oblik zapisa; racionalne vrednosti time nisu isključene.

Razlomljeni deo nastaje sabiranjem $2/7+3/49=(14+3)/49=17/49$. Iako su u osnovi $7$ dovoljne dve razlomljene cifre, broj cifara u drugoj osnovi nije određen tim brojem mesta. Konverzija čuva vrednost, a konačnost zapisa zavisi od osnove.
